\documentclass[11pt,a4paper]{article}

\usepackage[utf8]{inputenc}
\usepackage[T1]{fontenc}
\usepackage[margin=1in]{geometry}
\usepackage{amsmath,amssymb}
\usepackage{graphicx}
\usepackage{booktabs}
\usepackage{longtable}
\usepackage{enumitem}
\usepackage{xcolor}
\usepackage[style=authoryear,backend=biber,maxcitenames=2,maxbibnames=10,sorting=nyt]{biblatex}
\addbibresource{references.bib}
\usepackage{hyperref}
\hypersetup{colorlinks=true,linkcolor=blue!50!black,urlcolor=blue!50!black,citecolor=blue!50!black}

% --- Editorial note command: highlights TODOs / decisions for the senior researcher ---
\newcommand{\note}[1]{\textcolor{red}{\textbf{[NOTE: #1]}}}
\newcommand{\todo}[1]{\textcolor{blue!60!black}{\textbf{[TODO: #1]}}}

% --- Per-paper notes block -----------------------------------------------
% Usage: \begin{paper}{bibkey}{literature/bibkey.pdf}
%          \item ...free-form notes...
%        \end{paper}
\newenvironment{paper}[2]{%
  \subsection{\citetitle{#1}}%
  \label{paper:#1}%
  \noindent\textbf{\citeauthor{#1} (\citeyear{#1})} --- \citefield{#1}{journaltitle} \parencite{#1}\\
  \noindent\small\texttt{\detokenize{#2}}\normalsize
  \begin{itemize}[leftmargin=1.4em]
}{%
  \end{itemize}
  \paragraph{Notes}\todo{add notes}
  \medskip
}

% --- Citable fact / statement ----------------------------------------------
% Usage: \fact{bibkey}{Statement written as a manuscript-ready sentence.}
% Renders the sentence with an author-year citation, and prints the raw
% \parencite{bibkey} snippet so the line can be copy-pasted into a manuscript.
\newcommand{\fact}[2]{%
  \item #2~\parencite{#1}
    \\[-2pt]{\footnotesize\color{gray}\texttt{\textbackslash parencite\{#1\}}}%
}

\title{\textbf{Effect of Auditory Environment on Mental Workload \& Task Performance: 
    non-insvie mobile neuroimaging enables this study in naturalistic environments}\\[4pt]
\large Functional Hemispheric Asymmetry and Multisensory
Neurofeedback Training to Improve Brain Health\\[6pt]
\normalsize \emph{A Scoping Review Protocol and Research Framework}}

\author{[Senior Researcher Name]\\ \small [Affiliation]}
\date{Draft --- \today}

\begin{document}
\maketitle

\begin{abstract}
\noindent \todo{Write structured abstract last (Background / Objective / Methods /
Anticipated results). Keep to $\sim$250 words; follow the PRISMA-ScR abstract checklist.}
\smallskip

\noindent
\end{abstract}

\tableofcontents
\bigskip \hrule
\bigskip

% ============================================================
\section{Introduction and Rationale}
\label{sec:intro}

\subsection{Problem statement}
\begin{itemize}[leftmargin=1.4em]
    \item The brian is constantly processing the entire auditory environement, yet 
        attention and task performance is still achieved.
    \item We know that noisy enviroments are discrtacting and reduces perforamnce,
        but how and why?
    \item Gold-standard fMRI is itself a noisy machine requiring supnne
        subjects, thus mobile neurooimaging enables more naturalist study of
        these effects.
    \item EEG-fNIRS can enable this study. 
\end{itemize}


\section{Literature notes}
\label{sec:litnotes}
\todo{One block per PDF in \texttt{literature/}. Fill in the \emph{Notes} paragraph as you read; add \texttt{\textbackslash item}s for key findings, methods, sample, relevance.}

\begin{paper}{moore2017fatigue}{literature/moore2017fatigue.pdf}
  \item EEG/ERP (ERN, N1) + self-report + RT during 50-min auditory choice task; fatigue $\to$ reduced N1 (arousal) and decreased activation in dorsal attention network. Normal-hearing adults (n=19).
\end{paper}

\begin{paper}{alho2014auditorycortex}{literature/alho2014auditorycortex.pdf}
  \item Meta-analysis of fMRI studies: stimulus-dependent vs.\ attention-related modulation of auditory cortex. Useful background on where/how attention shapes auditory processing.
\end{paper}

\begin{paper}{yun2025anc}{literature/yun2025anc.pdf}
  \item fNIRS; active noise cancellation vs.\ noise during concentration tasks; neurophysiological effect of ANC in noisy environments. Very close to the target question.
    \item Auditory noisy environments is clearly linked to decreased task
        performance in humans, likely a consequence of increased stress and fatigue, yet neurophysiological evidence for the mechanics
        are not yet unassaiable. 
    \item Under noisy condictions the PFC has been shown to increase activity,
        potentially leading to ecxessive and unproductive mental workload.
    \item Inclusion criteria: pure-tone average (PTA) of less than 25 dB at 500,
        100, 2000 and 4000 Hz.
    \item Objective task perofmance + subjective questioneere
    \item HRFs via block-averaged fNIRS data -3 to +55 secnonds to task onset
    \item Brain activation quantified by $\beta$ via GLM with a modeled double gamma HRF and
    \item ANC ON vs ANC OFF in a two sample t-test. (+ point biserial
        correlation coefficent to gauge the time axis).
    \item channel signficance a Wilcoxon signed-rank test. Benjamin-Hochberg FDR
        correction, $\alpha = 0.05$
    \item ANC ON produced signficantly higher task performance than ANC OFF, but
        not signifiacnt differences in behvoural peforamnce such as reaction
        time. Indicating the congitivie load. 
    \item 3 of 4 channels with statistcally signficant difference between the
        two conditions were located in the right hemispehre (which corresponds
        to the literature about the right DLPFC in cogntive task and senstivity
        to cogntive load (MW)).
    \item Specifically, this study observed strong deactivation of HbO during
        the task, I beielive this should nuance a little. If the ANC ON had
        LOWER HBO than ANC OFF, what does that mean? 
    \item Subjective effort was signfivantly lower for ANC ON. 
    \item They conclude ANC OFF lead to unnescary consumption of neural
        resources, in light of the reduced task performance, this indicates ANC
        as effective to alleviate exxesvive cognitive load during audtiory noise. 
    \item However, this was a relativily simple task wihtout complexity, we know
        from the literature that neural effiancy and MW is senstivity to task difficulty. 
    \item THeir intepretation, that the stronger PFC deactvation during ANC OFF
        means that more of the brainresources was dedicated to the acutal task,
        rather than seperating speech from noise, thus refelcting the resuts.
        This opens a very intersting addition, temporal and pariental channels. 
    \item Further indicates right PFC as a senstivie area to cogntivie loads and
        MW. 
    \item Reccomends physiological measurements to avoid the variablity in
        subjective VAS listening effort measures. 
    \end{paper}

\begin{paper}{han2023entrainment}{literature/han2023entrainment.pdf}
  \item Combined acoustic + visual (acousto-optic) 40-Hz stimulation to drive coherent gamma entrainment. Relevant to the neurofeedback / multisensory-stimulation arm.
\end{paper}

\begin{paper}{hiwa2018fnirs}{literature/hiwa2018fnirs.pdf}
  \item fNIRS; auditory environments (e.g.\ silence/noise/music) vs.\ intellectual work performance; neural correlates in PFC.
  \item How music affects performance varies according to both the music and indiviudal.
    \item The efficancy of intellectual work is closely related to brain activation.
    \item Musical effect on intellectual work has been reported, but
        natrualistic noise environments are yet underreported. 
    \item Time of day + room humiditity and temprature reported, good. 
    \item 
  \end{paper}

\begin{paper}{gabbard2017fnirs}{literature/gabbard2017fnirs.pdf}
  \item fNIRS (PFC HbO) + eye tracking + RT/accuracy in simulated noisy ISR work; short-intermittent noise $\to$ larger PFC activation than long-intermittent noise.
  \item Audtiory noise increases cognitive load,, potentially reducing work
      performance by affecting attention, working memory, etc. 
    \item Again highlights that this doesnt not nessaercly hold for very simple
        tasks, touching neural effiancy theory. 
    \item MW can be asssed by PFC-based fNIRS. Stress can reduce PFC activation.
    \item Again $\alpha = 0.05$ 
    \item Their findings are not clear reagiding the "Noise leads to increased PFC
        activity" which I questioned immediatly and noted a questionmark in my notes. 
    \item Again task complecity is nuianced. It seems intermittent noise or
        bursts are more determimental than sustained noise.
    \item fNIRS is deemed a major canddicate for military applications to asses
        MW. 
    \item Specifcically this study found different types of noise affect the
        brain activation, supported by the Fan et al? 
    \item Rest-Memorize-Rest/Retain-Answer-Rest design, 8 digit short-term
        memory recall task. Again uses subjective VAS for pleasentneess,
        objective physiology would be benefincal here. 
    \item Used hitachy fnirs + POTATO for tomography analysis. Brain activation
        quantized via z-score. Functional connectivity by Pearsons into Fishers
        Z trasnform. 
    \item Again alpha 0.05. 
    \item Subjective pleasenett was signficantly higher in silence vs white
        noise. 
    \item Motor and visual cortiacal areas showed signficnalty stronger
        connecitvity in silence vs noise. 
    \item Strong indivudal differences led the analysis to consider 3 case:
        positive responsders, neutral and negative responders to the noise.
    \item For the positvie responsders (White group) no signficant brain
        activation difference between noise and silence was found. For the
        average and silence groups the brian activation differed signficantly
        between noise and silence in the right superior temporal gyrus,
        subcentral area, frontal pole, dorsolateral prefrontal cortex, and
        inferair frontal gyros. Again DLPFC seems to be a sensitive area to MW. 
    \item The silence group showed a signficnane difference in both
        temporal lobe with increased activation in silence. 
    \item The average group showed small difference in the pre-motor, sma areas
        between noise and silence
    \item The white group showed signcanfly increasd connectivity in the
        temporal areas during noise. The white group showed high
        interhemispheric connectivity during noise vs silence. 
    \item Average and Silence group showed signfiantly stronger PMA-SMA
        connectivity in silence versus noise (IMPORTANT), indicating motoric
        areas are downleguted during noise. 
    \item The silence group PMA-SMA autoconnetivity icnreased in noise, and
    \item Silece group showed higher DLPFC and pfc connectivity during silence vs
        noise (BIG)
    \item The White group shows complex habitation and compesation strategies
        enabling higher performance. 
    \item The brain activation between environments differed signficalty of rthe
        Silence group, not for the white group. 
    \item The average and silence groups showed decreased PFC-PMA-SMA connectivity
        during white noise versus silence (BIG). Here we see indicudal
        differences in neural effieancy or compesnatory behaviour, the average
        group's performance did not suffer from noise where the silence group
        did perform poorer during noise. Overall, the effect of noise can be
        individual and fNIRS may be an approriate tool to gauge how an
        individual responds to noise.m
    \end{paper}

\begin{paper}{ravi2021fnirs}{literature/ravi2021fnirs.pdf}
  \item Review-style paper on using fNIRS for cognitive workload assessment; surveys subjective vs.\ neuroimaging workload measures.
\end{paper}

\begin{paper}{pieper2021noise}{literature/pieper2021noise.pdf}
  \item EEG + subjective workload under environmental noise vs.\ noise-cancellation while working; EEG-based workload assessment.
\end{paper}

\begin{paper}{fan2022noise}{literature/fan2022noise.pdf}
  \item Noise exposure $\times$ mental workload; multiple physiological responses during task execution (check which modalities---EEG/ECG/EDA?).
\item Acuoustic noise defined is distrbuging, unpleasent or unwatned sound.
\item Acoustic noise exporuse is associated with negative health effects 
\item Physiology reacts signficantly to noise above 70 db. 
\item 35 to 65 db induces annoyces and subsequently can alter sleep patters. 
\item 66 to 85 db can cause signficant impairment including pshycological,
    neurological and hearing damage.
\item 85 to 115 db can cause signficant hearing damage and pshycological
    disturbance.
\item Mental workload (MW) in an optimal range improves human performance, where
    auditory noise is a signficnant factor in MW. THe mind wants an optimal
    range of MW, not under or overloaded. Flow state means an optimal amount of
    challenge to promote fun, learning and performance. 
\item MW can be measured in many objective and subjective methods, the article points to ECG (HRV), EEG and
    eyemovements.
\item Many have focused on dB meters and their effects, but rather we are
    interested in the characterizstics of the acoustic noise and what components
    consitutes the damaging or promotion of MW. 
\item A cogntiive task experiment was conducted with 3 noise levels N85-S1,
    N80-S1, and N75-S2, where N indictes decibel and S represents sharpness level. 
\item A NASA-authored cognitive task at different difficutlties were given.
\item ECG, EEG and eye tracking was monitored. I will only confeer ECG and EEG results.
\item Overall, the subjects used additional effort to complete the tasks under
    the noisy condictions and noise level contributes to MW. 
\item Noise intensity signficantly increased HRV.
\item Increased MW level decreased alpha power and increased theta power. 
\item Increased stress signfincaltly increased HRV and theta bandpower. -
\item Sharper noise decreased alpha power, i.e increase stress.
\item As MW incrases the body compoesnts by incrasing effort and stress, upto a
    given point where stress eentually impairs performance. 

\item Takeaway: Use ECG, EEG, fNIRS. Prioretize HRV, alpha \& theta band power,
    and PFC.
\end{paper}

\begin{paper}{causse2017workload}{literature/causse2017workload.pdf}
  \item fNIRS in PFC; mental workload and neural efficiency (piloting-type tasks); good methodological reference for fNIRS workload quantification.
  \item fNIRS is suitbale to measure MW and neural efficeancy in both labratroy
      and ecological settings.
    \item As the fNIRS signal did correlate strongly with percieved MW, but not
        actual task-perforamcen, indicates the already established existance of
        neural effieancy. 
    \item The neural efficancy index may be valid across indivudulas engaging in
        the same cognitive task, but breaks down when different cognitive domains 
        are at play.
    \item Increased task difficulity correletes to increased HbO and decreased
        HbR, indicating fNIRS are relevant for mental effort measurement. 
    \item However, predicting task performance from fNIRS alone is not feasable
        due to varaiblity in neural efficancy.
    \item See also \parencite{neubauer2009neuralefficiency} (neural efficiency review) and \S\ref{sec:facts}.
\end{paper}



\begin{paper}{rosenkranz2024eeg}{literature/rosenkranz2024eeg.pdf}
  \item Mobile EEG (ERPs, TRFs) during simulated laparoscopic surgery with OR soundscape; auditory work strain under low/high cognitive demand (n=22). Strong naturalistic-mobile-neuroimaging exemplar.
\end{paper}

\begin{paper}{noboa2025entrainment}{literature/noboa2025entrainment.pdf}
  \item Neural entrainment to beat + working memory predict sensorimotor synchronisation; relevant to auditory entrainment mechanisms.
\end{paper}

% NEW PAPER HERE TO ADD
% Cannon P, Ferrari E, Robie B, Michielli A, Minogue T, O'connor D. Quantification of cognitive load and usability in 3 percutaneous posterior spine surgical workflows using functional Near-Infrared Spectroscopy (fNIRS), Surgical Task Load Index (Surg-TLX) and modified System Usability Survey. Med Eng Phys. 2026 Aug 24. doi: 10.1088/1873-4030/ae9d83. Epub ahead of print. PMID: 42636870.
\item excessive cognitive load may negativly affect surgical performance
\item Subjective cogntivie load assment are valuable, but are retrospective and
    prone to influence by other factors. 
\item Thus, a multimodal approach may be beneficial to achieve relatime
    cognitive load assement. 
\item THe most common method is HRV objective measure, but fNIRS is proving itself 
\item This study aims not only to quantify and reduce mental fatigue and stress,
    but also guide the development and integration of technology into the
    surgical room where coupling task perforamnce, objective and subjective
    measures can guide proper implemetanions. 
\item They measure cogntivie load during robot-assisted (ARA), freehand (FHN) and
    floursopic-guide (FG) surgical proceedurs while measuring fNIRS. 
\item 6 surgeons, 
\item Assumes great PFC oxygenation means greater cogntive load, however we do
    know that there are some nuance here depending on task complexity and the
    indivudal 
\item fnirs metrics: Total peak oxygenation, duration-dependnt oxygenation,
    acerage oxygenation per second (FG vs ARA) and total peak haemoglobin
\item null-hypothesis is that the nmean of FHN eqauls ARA.
\item Result: the fNIRS metrics are signficant predictors of the srugical
    workflow and the null hypothesis is rejected. ARA signficanlty reduced
    cognitive load. 

% Fehring, D., Gaillard, A., Mazzoli, E. et al. Changes in prefrontal hemodynamics and mood states during screen use: a functional near-infrared spectroscopy study. Sci Rep 15, 28181 (2025). https://doi.org/10.1038/s41598-025-09360-w

% Dryden A, Allen HA, Henshaw H, Heinrich A. The Association Between Cognitive Performance and Speech-in-Noise Perception for Adult Listeners: A Systematic Literature Review and Meta-Analysis. Trends in Hearing. 2017;21. doi:10.1177/2331216517744675
  
% ============================================================
\section{Facts and statements}
\label{sec:facts}
\todo{Bank of citable, manuscript-ready sentences. One \texttt{\textbackslash fact\{bibkey\}\{sentence\}} per line; group by theme. Copy the sentence + the grey \texttt{\textbackslash parencite} snippet straight into the manuscript.}

\subsection{Rationale for mobile neuroimaging (vs.\ fMRI)}
\begin{itemize}[leftmargin=1.4em]
  \fact{tomasi2005fmrinoise}{The acoustic noise generated by the MRI scanner itself alters brain activation during working-memory tasks, confounding fMRI studies of auditory environment and cognition.}
\end{itemize}

\subsection{Noise, physiology and mental workload}
\begin{itemize}[leftmargin=1.4em]
  \fact{fan2022noise}{Acoustic noise is defined as disturbing, unpleasant or unwanted sound, and chronic noise exposure is associated with negative health effects.}
  \fact{fan2022noise}{Physiological responses to noise become significant above roughly 70~dB; 35--65~dB induces annoyance and can disrupt sleep, 66--85~dB can cause psychological, neurological and auditory impairment, and 85--115~dB can cause significant hearing damage.}
  \fact{fan2022noise}{Under noisy conditions participants exert additional effort to maintain task performance; noise intensity increases heart-rate-variability-based stress indices, and higher mental workload decreases EEG alpha power while increasing theta power.}
\end{itemize}

\subsection{Auditory environment and task performance}
\begin{itemize}[leftmargin=1.4em]
  \item \todo{add facts from \parencite{hiwa2018fnirs,gabbard2017fnirs,yun2025anc,pieper2021noise} as you read them.}
\end{itemize}

\subsection{Neural efficiency}
\begin{itemize}[leftmargin=1.4em]
  \fact{neubauer2009neuralefficiency}{The neural efficiency hypothesis that more able individuals show lower task-related brain activation holds primarily for tasks of low to moderate difficulty; in highly complex tasks, or in well-trained individuals, the relationship can invert such that increased activation accompanies better performance.}
  \fact{causse2017workload}{Prefrontal fNIRS signals correlate strongly with perceived mental workload but not with actual task performance, consistent with neural efficiency: individuals differ in how much activation they need to achieve the same outcome.}
  \fact{causse2017workload}{A neural efficiency index may be comparable across individuals performing the same cognitive task, but it does not transfer across different cognitive domains.}
  \fact{causse2017workload}{Increasing task difficulty produces increased prefrontal HbO and decreased HbR, supporting fNIRS as a measure of mental effort in both laboratory and ecological settings; however, predicting task performance from fNIRS alone is not feasible because of inter-individual variability in neural efficiency.}
\end{itemize}


% ============================================================

\nocite{arksey2005scoping,levac2010scoping,tricco2018prismascr}
\printbibliography[heading=bibintoc,title={References}]
\note{Scoping-review methodology refs (Arksey \& O'Malley 2005; Levac et al.\ 2010; Tricco et al.\ 2018) are in \texttt{references.bib} and forced into the list via \texttt{\textbackslash nocite}. Still to add: fNIRS-BCI review; neurofeedback review; sonification of neural signals; HAROLD / ageing lateralization; MCI \& neurodegeneration monitoring.}

\end{document}
